\chapter{کارهای مرتبط و پژوهش‌های پیشین}
\label{chap:related-work}

\section{مقدمه}
\label{sec:related-work-introduction}

پایش و تحلیل عملکرد برنامه‌های جاوا را می‌توان از دیدگاه‌های متفاوتی
انجام داد. برخی ابزارها بر رخدادهای داخلی ماشین مجازی تمرکز دارند، برخی
با نمونه‌برداری از پشته اجرا نقاط پرمصرف را شناسایی می‌کنند و گروهی دیگر
با استفاده از امکانات سیستم‌عامل، رفتار پردازنده، حافظه، فایل و شبکه را
در سطح پایین‌تر مشاهده می‌کنند. در کنار این ابزارها، چارچوب‌های
مشاهده‌پذیری معمولاً بر تولید و انتقال معیار، گزارش و رد توزیع‌شده تمرکز
دارند و سامانه‌های تشخیص ناهنجاری نیز داده‌های حاصل را برای یافتن رفتارهای
غیرعادی تحلیل می‌کنند.

هدف این فصل، معرفی فناوری‌های پایه نیست؛ این موضوع در فصل قبل انجام شده
است. در این فصل، راهکارهای موجود از نظر مسئله‌ای که حل می‌کنند، سطح
مشاهده، شیوه جمع‌آوری داده، مناسب‌بودن برای استفاده پیوسته، قابلیت پردازش
بلادرنگ و امکان تشخیص خودکار ناهنجاری مقایسه می‌شوند. این مقایسه کمک
می‌کند جایگاه سامانه پیاده‌سازی‌شده در این پایان‌نامه و تفاوت آن با
ابزارهای موجود روشن شود.

\section{ابزارهای مبتنی بر رخدادهای ماشین مجازی جاوا}
\label{sec:related-jvm-event-tools}

\subsection{\lr{Java Flight Recorder} و \lr{Java Mission Control}}

\lr{Java Flight Recorder} چارچوب ثبت رخداد در سکوی جاوا است و اطلاعاتی
درباره نخ‌ها، تخصیص حافظه، جمع‌آوری زباله، قفل‌ها، عملیات ورودی و خروجی و
فعالیت‌های داخلی ماشین مجازی تولید می‌کند. \lr{Java Mission Control}
نیز ابزار تحلیل و نمایش داده‌های ثبت‌شده توسط \lr{JFR} است. این دو ابزار
در کنار یکدیگر امکان ثبت پیوسته اطلاعات کم‌سطح و تحلیل پس از وقوع رخداد
را فراهم می‌کنند \cite{jfrdocs,jmcdocs,openjdkjmc}.

مزیت اصلی این رویکرد، یکپارچگی مستقیم با ماشین مجازی و دسترسی به رخدادهای
ساخت‌یافته‌ای است که ابزارهای عمومی سیستم‌عامل از آن‌ها آگاه نیستند.
همچنین، تنظیم نوع رخدادها و آستانه ثبت باعث می‌شود میان جزئیات داده و
سربار تعادل ایجاد شود. با این حال، \lr{JMC} عمدتاً یک ابزار تعاملی برای
تحلیل ضبط‌ها است و به‌تنهایی یک خط لوله مستقل برای انتقال رخدادها،
تجمیع جریانی و تشخیص ناهنجاری‌های چندمنبعی فراهم نمی‌کند.

سامانه این پایان‌نامه نیز از \lr{JFR} استفاده می‌کند، اما داده را تنها
برای مشاهده دستی در یک ابزار گرافیکی نگه نمی‌دارد. رخدادها در زمان اجرا
دریافت، به پیام ساخت‌یافته تبدیل و وارد مسیر جریانی می‌شوند. بنابراین،
\lr{JFR} در این سامانه نقش منبع داده را دارد و پردازش، نگه‌داری و تحلیل
در مؤلفه‌های مستقل انجام می‌شوند.

\subsection{\lr{VisualVM}}

\lr{VisualVM} یک ابزار دیداری برای عیب‌یابی و پروفایل‌سازی برنامه‌های
جاوا است. این ابزار قابلیت‌هایی مانند مشاهده پردازنده و حافظه، نمایش
نخ‌ها، تهیه \lr{heap dump} و \lr{thread dump} و پروفایل‌سازی بر اساس
نمونه‌برداری یا ابزارگذاری را ارائه می‌دهد \cite{visualvmdocs}. مزیت
آن، سهولت استفاده و فراهم‌کردن چند قابلیت تشخیصی در یک رابط واحد است.
برای توسعه‌دهنده‌ای که قصد بررسی تعاملی یک فرایند محلی یا از راه دور را
دارد، این ابزار می‌تواند اطلاعات اولیه و مناسبی فراهم کند.

محدودیت اصلی \lr{VisualVM} نسبت به مسئله این پژوهش، ماهیت ابزارمحور و
تعاملی آن است. این ابزار برای تحلیل دستی و مشاهده مستقیم مناسب است، اما
یک معماری توزیع‌شده برای انتقال پیوسته رخدادها، نگه‌داری آن‌ها در یک
سامانه پیام‌رسان و اجرای قواعد مستقل تشخیص ناهنجاری ارائه نمی‌دهد.
همچنین، تمرکز آن بر سطح جاوا است و مشاهده جامع رخدادهای سطح هسته نیازمند
ابزارهای تکمیلی خواهد بود.

\section{پروفایل‌سازهای نمونه‌بردار}
\label{sec:related-sampling-profilers}

\subsection{\lr{async-profiler}}

\lr{async-profiler} یک پروفایل‌ساز نمونه‌بردار کم‌سربار برای برنامه‌های
مبتنی بر \lr{HotSpot} است. این ابزار برای کاهش مسئله سوگیری نقاط امن%
\LTRfootnote{Safepoint Bias}
طراحی شده و می‌تواند نمونه‌های پردازنده، تخصیص حافظه، قفل‌ها و برخی
رخدادهای سخت‌افزاری یا سیستم‌عامل را جمع‌آوری کند. همچنین، پشته‌های جاوا،
بومی و هسته را در یک نمای مشترک نشان می‌دهد \cite{asyncprofiler}.

\lr{async-profiler} برای یافتن مسیرهای داغ، تحلیل مصرف پردازنده و تولید
نمودار شعله‌ای بسیار مناسب است. مشاهده هم‌زمان قاب‌های جاوا و بومی یکی از
مزیت‌های مهم آن نسبت به پروفایل‌سازهایی است که فقط پشته جاوا را ثبت
می‌کنند. با این حال، خروجی اصلی آن پروفایل آماری است و هدفش ساخت یک
سامانه عمومی انتقال رخداد، تجمیع چندنوع معیار در پنجره‌های زمانی و اجرای
آشکارسازهای ناهنجاری نیست.

پروفایل نمونه‌ای همچنین ماهیتی احتمالی دارد. رخدادهای کوتاه یا مسیرهایی
که میان دو نمونه اجرا می‌شوند ممکن است در خروجی دیده نشوند. پژوهش‌های
مقایسه‌ای درباره پروفایل‌سازهای جاوا نیز نشان داده‌اند که ابزارهای مختلف
ممکن است در تشخیص متدهای داغ یا میزان سربار با یکدیگر اختلاف داشته باشند
\cite{donttrustprofiler}. بنابراین، پروفایل نمونه‌ای باید با شناخت
محدودیت‌های روش و هدف تحلیل استفاده شود.

\subsection{چارچوب‌های ابزارگذاری ماشین مجازی}

رابط‌های ابزارگذاری ماشین مجازی مانند \lr{JVMTI} امکان دریافت اطلاعات
جزئی درباره اجرای برنامه، نخ‌ها، حافظه و رخدادهای ماشین مجازی را فراهم
می‌کنند. این رابط‌ها پایه ساخت بسیاری از عامل‌ها و پروفایل‌سازهای جاوا
هستند و می‌توانند برای ثبت رخدادهایی مانند ورود و خروج متد یا تخصیص حافظه
به کار روند \cite{jvmtipaper}.

مزیت ابزارگذاری، امکان کنترل دقیق نقاط جمع‌آوری و دریافت داده‌های سطح
متد است. در مقابل، فعال‌سازی رخدادهای پرتکرار مانند ورود و خروج همه متدها
می‌تواند سربار قابل توجهی ایجاد کند. همچنین، توسعه عامل بومی و مدیریت
سازگاری با نسخه‌های مختلف ماشین مجازی پیچیدگی بیشتری نسبت به مصرف
رخدادهای آماده \lr{JFR} دارد. سامانه این پژوهش برای کاهش نیاز به
ابزارگذاری فراگیر، از رخدادهای \lr{JFR} به‌عنوان منبع اصلی سطح ماشین
مجازی استفاده می‌کند.

\section{ابزارهای مشاهده‌پذیری مبتنی بر \lr{eBPF}}
\label{sec:related-ebpf-tools}

\subsection{\lr{BCC}}

\lr{BCC} یا مجموعه کامپایلر \lr{BPF}، چارچوبی برای ساخت ابزارهای ردیابی
پویا در لینوکس است. این مجموعه ابزارهای آماده‌ای برای بررسی فرایند،
پردازنده، فایل، شبکه و حافظه ارائه می‌کند و امکان نوشتن برنامه‌های سفارشی
برای اتصال به نقاط ردیابی هسته یا فضای کاربر را نیز فراهم می‌سازد
\cite{bccdocs,bccprimer}.

مزیت \lr{BCC} انعطاف‌پذیری زیاد و دسترسی به رخدادهای سطح پایین است. با
این حال، بسیاری از ابزارهای آن برای اجرای تعاملی توسط متخصص عملکرد طراحی
شده‌اند و خروجی هر ابزار قالب و هدف خاص خود را دارد. ایجاد یک سامانه
پایدار که رخدادهای منتخب را پیوسته دریافت، نرمال‌سازی و به زیرساخت
پردازش جریانی ارسال کند، همچنان نیازمند طراحی و پیاده‌سازی اضافی است.

در پروژه این پایان‌نامه، نقش مشابهی با برنامه‌های اختصاصی \lr{eBPF} و
یک فرستنده فضای کاربر پیاده‌سازی شده است. تفاوت اصلی آن است که رخدادهای
انتخاب‌شده مستقیماً وارد مدل داده و خط لوله مشترک سامانه می‌شوند و در
کنار رخدادهای \lr{JFR} قابل تجمیع و تحلیل هستند.

\subsection{\lr{bpftrace}}

\lr{bpftrace} یک زبان سطح بالا برای ردیابی پویا در لینوکس است که بر
\lr{eBPF} تکیه دارد و نوشتن ردیاب‌های کوتاه برای نقاطی مانند
\lr{kprobe}، \lr{uprobe} و \lr{tracepoint} را ساده می‌کند
\cite{bpftracedocs}. این ابزار برای بررسی سریع فرضیه‌ها، تحلیل رخدادهای
موقت و ساخت نمونه اولیه بسیار مناسب است.

با وجود سهولت استفاده، اسکریپت‌های \lr{bpftrace} معمولاً برای تحلیل
تعاملی و مسئله‌محور نوشته می‌شوند. تبدیل آن‌ها به یک مؤلفه پایدار با مدل
پیام نسخه‌پذیر، کنترل ازدحام، انتقال شبکه‌ای و بازیابی خطا نیازمند
طراحی‌های تکمیلی است. پروژه حاضر به جای تکیه بر اسکریپت‌های تعاملی، مسیر
مشخصی برای ارسال رخداد از هسته به فضای کاربر و سپس انتشار در
\lr{Kafka} ایجاد کرده است.

\subsection{\lr{Pixie}}

\lr{Pixie} یک سکوی مشاهده‌پذیری برای محیط‌های \lr{Kubernetes} است که با
استفاده از \lr{eBPF} بخشی از داده‌های تله‌متری را بدون ابزارگذاری دستی
برنامه جمع‌آوری می‌کند. این ابزار می‌تواند اطلاعاتی مانند ترافیک
برنامه‌ها، وضعیت سرویس‌ها، مصرف منابع و نمودارهای شعله‌ای را نمایش دهد
\cite{pixiedocs,pixiearchitecture}.

ویژگی مهم \lr{Pixie} خودکارسازی جمع‌آوری تله‌متری در سطح خوشه و ارائه
قابلیت پرس‌وجو و نمایش نزدیک به زمان واقعی است. در مقابل، دامنه اصلی آن
محیط \lr{Kubernetes} و مشاهده تعامل سرویس‌ها و پادها است. هدف آن
پروفایل‌سازی اختصاصی یک فرایند جاوا با ترکیب رخدادهای ساخت‌یافته
\lr{JFR} و مجموعه‌ای از آشکارسازهای ناهنجاری ویژه منابع نیست. ازاین‌رو،
اگرچه در استفاده از \lr{eBPF} و پردازش نزدیک به زمان واقعی با این پژوهش
اشتراک دارد، سطح انتزاع و مسئله هدف متفاوت است.

\section{چارچوب‌های عمومی مشاهده‌پذیری}
\label{sec:related-observability-frameworks}

\subsection{\lr{OpenTelemetry}}

\lr{OpenTelemetry} یک چارچوب مستقل از فروشنده برای تولید، جمع‌آوری و
ارسال ردها، معیارها و گزارش‌ها است. در اکوسیستم جاوا، این چارچوب
ابزارگذاری دستی و عامل جاوا برای ابزارگذاری بدون تغییر مستقیم کد را
فراهم می‌کند \cite{opentelemetrydocs,oteljavaagent}. عامل جاوا با
ابزارگذاری پویا می‌تواند تله‌متری مربوط به درخواست‌های ورودی، فراخوانی‌های
خروجی، پایگاه داده و کتابخانه‌های متداول را ثبت کند.

\lr{OpenTelemetry} برای مشاهده‌پذیری سرویس‌ها، ردیابی درخواست‌های
توزیع‌شده و استانداردسازی انتقال تله‌متری بسیار مناسب است. با این حال،
تمرکز آن با پروفایل‌سازی رخدادمحور \lr{JVM} و ردیابی سطح هسته یکسان
نیست. داده‌های سطح متد، رخدادهای داخلی جمع‌آوری زباله یا رخدادهای خاص
هسته به‌طور پیش‌فرض مسئله اصلی این چارچوب نیستند. همچنین،
\lr{OpenTelemetry} خود یک سامانه تشخیص ناهنجاری نیست و معمولاً داده را
برای ذخیره‌سازها و سامانه‌های تحلیل بیرونی صادر می‌کند.

سامانه این پژوهش از نظر جداسازی تولید تله‌متری و تحلیل با رویکرد عمومی
مشاهده‌پذیری شباهت دارد، اما منبع داده و هدف آن تخصصی‌تر است. رخدادهای
\lr{JFR} و \lr{eBPF} به معیارهای عملکردی پنجره‌ای تبدیل شده و سپس توسط
آشکارسازهای اختصاصی تحلیل می‌شوند.

\subsection{\lr{Prometheus}}

\lr{Prometheus} یک سامانه پایش و پایگاه داده سری زمانی است که معیارها را
معمولاً با الگوی کششی از نقاط پایانی جمع‌آوری می‌کند و امکان پرس‌وجو و
تعریف قواعد هشدار را فراهم می‌سازد \cite{prometheusdocs}. این ابزار برای
ثبت معیارهای عددی، نمایش روند و هشداردهی بر اساس عبارات مشخص بسیار
مناسب است.

تفاوت مهم معیارهای \lr{Prometheus} با رخدادهای خام این پژوهش در سطح
جزئیات است. \lr{Prometheus} معمولاً مقادیر تجمیع‌شده را جمع‌آوری می‌کند،
در حالی که خط لوله این پژوهش ابتدا رخدادهای ساخت‌یافته را دریافت و سپس
بر اساس کلید و پنجره زمانی آن‌ها را به معیار تبدیل می‌کند. در پروژه،
\lr{Prometheus} برای پایش وضعیت خود مؤلفه‌ها به کار می‌رود، اما مسیر
اصلی تحلیل ناهنجاری بر داده‌های جریانی \lr{Kafka} متکی است.

\section{روش‌های تشخیص ناهنجاری در داده‌های پایشی}
\label{sec:related-anomaly-detection}

تشخیص ناهنجاری حوزه گسترده‌ای است و روش‌های آن از قواعد ساده تا مدل‌های
یادگیری عمیق را در بر می‌گیرند. یک مرور شناخته‌شده، ناهنجاری‌ها را بر
اساس ماهیت داده و دسترسی به برچسب‌ها بررسی کرده و روش‌ها را در دسته‌هایی
مانند آماری، مبتنی بر فاصله، خوشه‌بندی و طبقه‌بندی قرار می‌دهد
\cite{chandolasurvey}. در داده‌های سری زمانی، ترتیب زمانی، روند،
فصل‌مندی و وابستگی میان نقاط اهمیت ویژه دارد.

روش آستانه‌ای یکی از ساده‌ترین رویکردها است. در این روش، مقدار معیار با
یک حد ثابت یا پویا مقایسه می‌شود. مزیت آن هزینه پایین و قابلیت توضیح
مستقیم است، اما انتخاب آستانه به محیط و بار کاری وابسته است. خط پایه
غلتان این محدودیت را تا حدی کاهش می‌دهد و مقدار جاری را با رفتار چند
بازه قبلی مقایسه می‌کند. تحلیل روند نیز برای مشکلاتی مانند رشد تدریجی
حافظه مناسب است که الزاماً با یک جهش ناگهانی ظاهر نمی‌شوند.

روش‌های یادگیری ماشین و یادگیری عمیق می‌توانند الگوهای پیچیده‌تر و روابط
چندمتغیره را بیاموزند. مرورهای جدید در تشخیص ناهنجاری سری زمانی نشان
می‌دهند که مدل‌های عمیق توانایی نمایش الگوهای غیرخطی را دارند، اما به
داده، تنظیمات، هزینه محاسباتی و سازوکار ارزیابی دقیق نیازمندند
\cite{deepanomalysurvey}. همچنین، قابلیت توضیح نتیجه در بسیاری از
مدل‌های پیچیده کمتر از قواعد مستقیم است.

در سامانه این پایان‌نامه، روش‌های آستانه‌ای، خط پایه غلتان و تشخیص روند
انتخاب شده‌اند. این انتخاب با هدف پروژه، یعنی ایجاد یک نمونه عملی
بلادرنگ و قابل توضیح، سازگار است. روش مورد استفاده ادعای حل عمومی مسئله
تشخیص ناهنجاری را ندارد، اما امکان مشخص‌کردن معیار، مقدار، آستانه و بازه
وقوع ناهنجاری را فراهم می‌کند. این قابلیت برای ارزیابی سناریوهای
کنترل‌شده و تحلیل مهندسی خروجی اهمیت دارد.

\section{مقایسه راهکارهای موجود}
\label{sec:related-work-comparison}

جدول~\ref{tab:related-work-comparison} ابزارها و رویکردهای بررسی‌شده را
از نظر سطح مشاهده و قابلیت‌های اصلی مقایسه می‌کند. منظور از «تشخیص
ناهنجاری» در این جدول، وجود سازوکار داخلی برای تحلیل خودکار رفتار
غیرعادی است، نه صرفاً امکان مشاهده داده توسط کاربر.

\begin{table}[H]
    \centering
    \caption{مقایسه ابزارها و رویکردهای مرتبط}
    \label{tab:related-work-comparison}
    \small
    \begin{tabular}{|p{2.4cm}|p{2.5cm}|p{3cm}|p{2.5cm}|p{2.4cm}|}
        \hline
        \textbf{ابزار یا رویکرد} &
        \textbf{سطح مشاهده} &
        \textbf{شیوه اصلی استفاده} &
        \textbf{پردازش پیوسته} &
        \textbf{تشخیص ناهنجاری داخلی} \\
        \hline

        \lr{JFR/JMC} &
        ماشین مجازی جاوا &
        ثبت رخداد و تحلیل تعاملی &
        ثبت پیوسته ممکن است، اما تحلیل غالباً تعاملی است &
        ندارد \\
        \hline

        \lr{VisualVM} &
        فرایند و ماشین مجازی جاوا &
        مشاهده و پروفایل‌سازی تعاملی &
        محدود به نشست ابزار &
        ندارد \\
        \hline

        \lr{async-profiler} &
        پشته جاوا، بومی و هسته &
        نمونه‌برداری و تولید پروفایل &
        قابل اجرای پیوسته با مدیریت خروجی &
        ندارد \\
        \hline

        \lr{BCC/bpftrace} &
        هسته و فضای کاربر لینوکس &
        ردیابی پویا و تحلیل مسئله‌محور &
        ممکن است، اما نیازمند ساخت خط لوله اختصاصی است &
        ندارد \\
        \hline

        \lr{Pixie} &
        خوشه \lr{Kubernetes} و ترافیک برنامه‌ها &
        مشاهده‌پذیری خودکار مبتنی بر \lr{eBPF} &
        دارد &
        محدود و وابسته به قابلیت‌های سکوی تحلیل \\
        \hline

        \lr{OpenTelemetry} &
        رد، معیار و گزارش برنامه و سرویس &
        ابزارگذاری و صدور تله‌متری &
        دارد &
        ندارد؛ به سامانه تحلیل بیرونی نیاز دارد \\
        \hline

        سامانه این پژوهش &
        \lr{JVM} و هسته لینوکس &
        دریافت رخداد، انتقال، تجمیع و تحلیل &
        دارد &
        قواعد آستانه‌ای، خط پایه و روند \\
        \hline
    \end{tabular}
\end{table}

این مقایسه نشان می‌دهد که ابزارهای تخصصی معمولاً در حوزه اصلی خود اطلاعات
عمیق‌تری ارائه می‌کنند. برای مثال، \lr{async-profiler} برای ساخت
پروفایل پردازنده و \lr{Pixie} برای مشاهده‌پذیری محیط
\lr{Kubernetes} قابلیت‌هایی فراتر از نیازهای محدود پروژه حاضر دارند.
بنابراین، هدف سامانه این پایان‌نامه جایگزینی این ابزارها نیست؛ بلکه
یکپارچه‌کردن چند مرحله مشخص از جمع‌آوری رخداد تا گزارش ناهنجاری در یک
نمونه قابل اجرا است.

\section{شکاف موجود و جایگاه پژوهش}
\label{sec:research-gap}

بررسی کارهای مرتبط چند شکاف را نشان می‌دهد. نخست، بسیاری از ابزارهای
پروفایل‌سازی جاوا بر یک نشست تحلیل یا یک خروجی پروفایل تمرکز دارند و
مسئولیت انتقال، ذخیره و تحلیل پیوسته داده را بر عهده کاربر می‌گذارند.
دوم، ابزارهای \lr{eBPF} دید عمیقی از سیستم‌عامل فراهم می‌کنند، اما از
مفاهیم داخلی ماشین مجازی مانند رخدادهای جمع‌آوری زباله، تخصیص شیء و
قفل‌های جاوا آگاه نیستند. سوم، چارچوب‌های عمومی مشاهده‌پذیری بر
استانداردسازی معیار، گزارش و رد تمرکز دارند و الزاماً داده‌های
پروفایل‌سازی یا رخدادهای سطح هسته را به یک مدل تحلیلی مشترک تبدیل
نمی‌کنند.

همچنین، در بسیاری از ابزارها تحلیل نهایی به کاربر واگذار می‌شود.
نمودار، پروفایل یا رخداد خام اطلاعات ارزشمندی ارائه می‌دهد، اما برای
شناسایی پیوسته رفتار غیرعادی لازم است قواعد یا مدل‌های تحلیل روی جریان
داده اجرا شوند. اتصال منبع داده به یک سامانه پردازش جریانی و سپس اجرای
آشکارسازهای مستقل، این فاصله میان مشاهده و تصمیم‌گیری اولیه را کاهش
می‌دهد.

جایگاه پژوهش حاضر در ترکیب این اجزا است: رخدادهای ساخت‌یافته
\lr{JFR} و رخدادهای منتخب \lr{eBPF} از دو مسیر مستقل جمع‌آوری می‌شوند؛
پس از نرمال‌سازی وارد \lr{Kafka} می‌گردند؛ با \lr{Kafka Streams} در
پنجره‌های زمانی تجمیع می‌شوند؛ و خروجی توسط آشکارسازهای قابل توضیح تحلیل
می‌شود. این ترکیب یک الگوریتم بنیادی جدید برای پروفایل‌سازی یا تشخیص
ناهنجاری ارائه نمی‌کند، اما یک نمونه انتهابه‌انتها و توسعه‌پذیر برای
پایش چندلایه برنامه‌های جاوا فراهم می‌سازد.

\section{جمع‌بندی}
\label{sec:related-work-summary}

در این فصل، ابزارهای مبتنی بر رخدادهای ماشین مجازی، پروفایل‌سازهای
نمونه‌بردار، ابزارهای ردیابی مبتنی بر \lr{eBPF}، چارچوب‌های عمومی
مشاهده‌پذیری و روش‌های تشخیص ناهنجاری بررسی شدند. ابزارهایی مانند
\lr{JMC}، \lr{VisualVM} و \lr{async-profiler} دید مناسبی از رفتار
برنامه جاوا فراهم می‌کنند، در حالی که \lr{BCC}، \lr{bpftrace} و
\lr{Pixie} از امکانات \lr{eBPF} برای مشاهده سیستم‌عامل و محیط اجرا
استفاده می‌کنند. \lr{OpenTelemetry} و \lr{Prometheus} نیز انتقال
تله‌متری و پایش معیارمحور را تسهیل می‌کنند.

هیچ‌یک از این ابزارها به‌تنهایی دقیقاً همان زنجیره‌ای را که در این
پایان‌نامه پیاده‌سازی شده است پوشش نمی‌دهد. سامانه حاضر رخدادهای
\lr{JVM} و هسته را در یک معماری جریانی قرار می‌دهد و داده تجمیع‌شده را
با قواعد قابل توضیح تحلیل می‌کند. فصل بعد با تکیه بر مبانی نظری و
مقایسه‌های انجام‌شده، معماری این سامانه و جریان داده میان مؤلفه‌های آن
را تشریح می‌کند.
